\chapter{Introduction}\label{ch:introduction}

\section{Motivation and task}
Automated news routing systems must sort incoming feeds accurately even when rushed copy editing, optical character recognition glitches, or hurried keyboard slips introduce small spelling errors. A human reader usually infers a slightly garbled word almost effortlessly from its outer letters and sentence context. A machine classifier enjoys no such leeway unless its internal representation preserves partial word similarity or absorbs missing tokens gracefully. Once a single adjacent letter swap turns a domain keyword into a string absent from training data, discrete lookup tables drop that clue immediately, and downstream phrase detectors or linear sums can shift toward the wrong category.

We investigate how strongly this degradation hits across four-topic news classification using the AG News corpus \citep{zhang2015,agmirror}. Each classifier reads one article string formed from the story headline and short summary and assigns it to one of four balanced categories: \emph{World}, \emph{Sports}, \emph{Business}, or \emph{Sci/Tech}. To model realistic typing noise without changing article length, word boundaries, or human readability, we focus on adjacent interior letter transpositions, swapping two unequal neighboring characters strictly between the first and last letters of eligible words (turning \texttt{consumer} into \texttt{cosnumer}, for instance). Holding outer anchor letters fixed keeps each word recognizable to a person while breaking exact string matches for any model built on whole words. Applying these edits deterministically per article guarantees that every classifier faces identical clean and noisy test inputs.

\section{Research questions}
Because model behavior depends heavily on training volume, feature design, and compute budget, our primary study fixes a clean baseline budget of 20,000 labeled training articles (5,000 per topic) alongside an untouched 4,000-article validation split and a 4,000-article confirmation split separated prior to final testing. Evaluating neural and linear architectures on equal labeled data and standard CPU hardware prevents extra external pretraining or massive label budgets from masking real inductive-bias differences. Around that shared experimental core, we test two central questions:
\begin{enumerate}
\item \emph{How quickly does accuracy erode under increasing interior swap rates, and how do word-level and subword-aware convolutional networks compare against a sparse linear baseline?} Measuring clean-trained models across 0\%, 5\%, 10\%, and 20\% word-corruption rates shows whether ordered phrase detectors retain context better than global word checklists when individual tokens break.
\item \emph{Does replacing half of the clean training articles with 10\%-swapped copies improve noisy test accuracy without sacrificing clean-text performance?} Holding the total training count fixed at 20,000 stories isolates whether training-time typo exposure builds robustness on its own rather than simply benefiting from extra rows.
\end{enumerate}

To make sure answers to those questions reflect genuine algorithmic differences rather than lucky random seeds or single test draws, our protocol tracks three seeded runs per primary configuration, checks two additional random 20\% test draws, verifies topic-balanced behavior alongside overall accuracy, and pairs clean and noisy predictions article by article. We also separate schedule length from subword modeling using a matched 16-epoch word-only neural control, ensuring any subword advantage is not merely an artifact of longer optimization.

\section{Approach and findings}
Our focused comparison evaluates three scratch-trained classifiers covering the main architectural trade-off between global unordered vocabulary scoring and ordered local phrase detection:
\begin{itemize}
\item \textit{Word TF-IDF LR} converts each normalized article into a high-dimensional sparse vector of single words and adjacent word pairs, weighting rare topical vocabulary more heavily than ubiquitous filler terms and summing learned class weights across the whole document.
\item \textit{Word CNN} preserves token order by embedding recognized vocabulary words into learned dense vectors, routing unfamiliar words and novel typos to one shared fallback vector, sliding phrase filters of three to five words over the sequence, and keeping each filter's single strongest match across valid positions.
\item \textit{Subword CNN} extends \textit{Word CNN} with learned vectors for hashed three-to-five-character fragments averaged into each word representation so an unseen misspelling still carries intact prefix, stem, and suffix signals into those phrase filters.
\end{itemize}

Before locking those three models into multi-run testing, we explored nine linear, perturbation-invariant, pretrained, and neural candidates on the shared validation split (Chapter~\ref{ch:experiments}), confirming both where scratch models sit relative to frozen sentence embeddings and why no single clean-trained variant escapes noise completely. Across all held-out evaluations, \textit{Word TF-IDF LR} dominates on both clean accuracy (90.05\%) and typo tolerance (89.14\% at 20\% swaps) while tuning in a median of 13 seconds on CPU. Both scratch-trained CNNs drop more than twice as sharply under 20\% swaps and require over an order of magnitude more training time. Adding training-time swaps yields a narrow gain of +0.46 percentage points at 20\% swaps for \textit{Subword CNN}, practically zero movement for \textit{Word TF-IDF LR}, and no consistent gain for \textit{Word CNN}.

The rest of this report develops these results step by step. Chapter~\ref{ch:background} explains the mechanics of sparse linear checklists, convolutional phrase detectors, character hashing, and human typo tolerance. Chapter~\ref{ch:method} details dataset auditing, deterministic corruption, neural padding masks, and formal training objectives. Chapter~\ref{ch:experiments} reviews our nine exploratory model families, learning dynamics, uncertainty estimation, and CPU runtimes. Chapter~\ref{ch:results} presents multi-seed accuracy curves, augmentation confidence intervals, confirmation checks, topic confusions, and concrete article error breakdowns, before Chapters~\ref{ch:discussion} and~\ref{ch:conclusion} analyze structural causes, schedule effects, factorial extensions, and final takeaways.
